\documentclass[11pt]{exam}
\usepackage[a4paper,margin={.1\paperheight,.1\paperwidth},marginratio=1:1]{geometry}
\usepackage[french]{babel}
\usepackage[T1]{fontenc}

% Choix des options et des infos d’en-tête
\usepackage[cours,prof]{profnsi}
\profnsisetup{
	niveau=Première,
	matiere=NSI,
	titre=Représentation des données,
	chapitre=3
}


\begin{document}
	\creerDoc
	
	\partie{Introduction : Pourquoi représenter les données ?}
	
	\souspartie{Intérêt}
	
	Les ordinateurs ne manipulent pas directement des nombres, des textes ou des images comme nous les percevons : ils ne savent traiter que des \textbf{états physiques} simples.  
	Concrètement, dans la plupart des circuits électroniques, un transistor peut être dans deux états distincts :
	
	\begin{itemize}
		\item un courant passe (\textbf{état haut}) ;
		\item aucun courant ne passe (\textbf{état bas}).
	\end{itemize}
	
	Ces deux états sont naturellement représentés par deux symboles : \(0\) et \(1\).  
	Ainsi, toute information manipulée par un ordinateur doit être transformée en une suite de 0 et de 1 : c’est ce que l’on appelle la \textbf{représentation binaire}.
	
\souspartie{remarques}

	\remarque{Pense à convertir en binaire avant de simplifier.}
	
	\remarque[Astuce]{Commence par les puissances de 2 les plus grandes.}
	
\souspartie{lignes et applications}
	% Lignes
	\lignes{3}
	% Applications
	\application{Convertis 37 en binaire.}{4}
	
	\application[Astuce]{Explique la conversion.}{4}
	
	\application[Astuce]{Explique la conversion.}
	
\souspartie{code}

\begin{python}
	def somme(a, b):
		return a + b
\end{python}

\begin{sql}
	SELECT * FROM persons WHERE age >= 18;
\end{sql}

\begin{html}
	<!DOCTYPE html>
	<html>
	<head>
	<title>Titre de la page</title>
	</head>
	<body>
	<h1>Texte principal</h1>
	<p>Mon paragraphe</p>
	</body>
	</html>
\end{html}

\begin{bash}
	ls -la
	cd /etc/nginx
	sudo systemctl restart nginx
\end{bash}




Dans le texte : \code{Python}{def x():}

\souspartie{trou}

Dans cet exemple, le langage utilisé est \trou{Python}{1}

Les données transférées sur internet sont \trou{découpées en paquets. Les paquets sont expédiés de l'expéditeur vers le destinataire en passant par des routeurs aux n\oe uds du réseaux}{3} qui décident vers quel autre n\oe ud il doit les envoyer ; un routeur choisit parmi les différents chemins possibles en fonction de l'encombrement du réseau, des pannes, etc.

\souspartie{pointillés}

% pleine largeur du texte
\pointilles

% largeur personnalisée
\pointilles[6cm]

% dans une phrase
1+1 : \pointilles[4cm]   \(4 \times 4\) : \pointilles[2cm]

\souspartie{Opérations}

% Décimal
\begin{center}
\addition{8475}{6520}

% Binaire, avec et sans retenues
\addition[Base=bin]{1111}{111}

\addition[Base=bin,AffRetenues=false]{1111}{111}

% Hexa, sans '=' et sans espacement horizontal
\addition[Base=hex,AffEgal=false,Offset=0pt]{ABCD}{FE}

% Soustraction binaire
\soustraction[Base=bin]{10000}{111}

% Multiplication hexa avec décalage '.'
\multiplication[Base=hex]{ABCD}{FAFD}

\multiplication[Base=bin]{11}{11}
\end{center}

\souspartie{Divisions}

\begin{center}

\division{128}

\division[3]{98}

\end{center}

\souspartie{image}

\img{img/image.jpg}

\img[0.2]{img/image.jpg}

\souspartie{tableaux}

\tableau{Décimal|Binaire}{4|100 ; 5|101 ; 6|110}

\tableau[.5\linewidth]
{Tour|i|compteur}
{
	0|0|0 ;
	1|1|1 ;
	10|2|2
}

\souspartie{arbres}

\arbre{[A [B] [C [E] [F]] [D]]}

\arbre{%
	[/
	[bin]
	[users [louis] [jean]]
	[dev]
	]
}

\begin{center}
	\arbre{%
		[Projet
		[Spécifications]
		[Dev [Backend] [Frontend]]
		[Tests]
		[Livraison]
		]
	}
\end{center}



\souspartie{graphes}

\begin{GrapheTikz}
	%SOMMETS
	\GrphPlaceSommets{(1,4)/A (3,4)/B (4,3)/C (3,2)/D (2,1)/E (0,1)/F}
	%ARÊTES
	\GrphTraceAretes{A/B B/A B/D B/E C/A C/D D/A D/E E/F F/B}
\end{GrapheTikz}


\souspartie{Matrices d'adjacence}

\matrice{0 & 1 & 0 \\ 1 & 0 & 1 \\ 0 & 1 & 0}

\matriceadj[G]{0 & 1 & 0 \\ 1 & 0 & 1 \\ 0 & 1 & 0}


\souspartie{piles - files}


Voici une pile :  \\
\pile{[2, 4, 3, 9, 5]}\\

Voici une file :  \\
\file{[A, V, P, Z, H]}

\souspartie{Listes chaînées}

Voici ma liste chaînée : \\

\listechainee{2, 4, 3, 9}

\souspartie{QCM}

\begin{qcm}
	
	\begin{qcmquestion}{La dérivée de $x^2$ est :}
		\begin{qcmchoices}
			\qcmchoice{$x$}
			\qcmcorrectchoice{$2x$}
			\qcmchoice{$x^3$}
			\qcmchoice{$\dfrac{1}{x}$}
		\end{qcmchoices}
		\qcmsolution{On a $(x^2)'=2x$.}
	\end{qcmquestion}
	
	\begin{qcmquestion}{La dérivée de $x^2$ est :}
		\begin{qcmchoices}
			\qcmchoice{$x$}
			\qcmcorrectchoice{$2x$}
			\qcmchoice{$x^3$}
			\qcmchoice{$\dfrac{1}{x}$}
		\end{qcmchoices}
		\qcmsolution{On a $(x^2)'=2x$.}
	\end{qcmquestion}
	
	\begin{qcmfreequestion}{Donner la matrice d'adjacence}[\matriceadj[G]{0 & 1 & 0 \\ 1 & 0 & 1 \\ 0 & 1 & 0}][5]
	\end{qcmfreequestion}
	
\end{qcm}

\newpage


% pleine largeur du texte
\pointilles
\\

% largeur personnalisée
\pointilles[6cm]
\\

% dans une phrase
1+1 : \pointilles[4cm]   \(4 \times 4\) : \pointilles[2cm]
\\

Le langage \pointillestrou{html}[6cm]

\addition[Base=bin]{00000101}{11111011}


\end{document}
